\chapter{\texorpdfstring{LITERATURE REVIEW}{}} %upper case only

This chapter provides a review of the literature that informed my investigation into US immigration policies, Pakistani immigrants, Pakistani family cultural values and religious practices. The chapter also examines how the existing relevant literature shaped my thinking about the roles of retrospective storytelling in the communication of identity and the theoretical framework. Finally, I will introduce the research questions, which emerged from the existing scholarship.

\section{US Immigration Policies}

US immigration policies have always influenced the demographic composition and the factors that have driven each wave of immigrants to the US. From the 1920s to the Civil Rights Movement of the 1960s, the US implemented a national origins quota system, limiting the number of immigrants allowed entry. Furthermore, \citet{} Vesely et al. (2016) shed light on the National Origins Act of 1924 that restricted immigration visas to 2 percent of the total number of people of each nationality represented in the 1890 US census. This act excluded immigrants from Asia and imposed restrictions on immigrants from Eastern Europe and Africa. The McCarran–Walter Act of 1952 ended Asian exclusion and introduced a system of preferences based on skills, and the Immigration and Nationality Act of 1965 (Hart–Celler Act) furthered a family reunification program for Asian immigrants \citep{batalovaImmigrationDataMatters2020} (Vesely et al., 2016).

Pakistani Americans’ presence in the US became more visible after the Immigration and Nationality Act of 1965 that allowed skilled professionals, particularly doctors and engineers, to immigrate from South Asia following changes in the US immigration policy \citep{batalovaImmigrationDataMatters2020, ghaniPakistaniMigrantsUnited2016, PakistaniImmigrantsSocial}. Those immigrants, who were either unable to qualify for regulated professions or came on tourist or student visas, and often reinvented themselves as small business owners, frequently operating convenience stores or gas stations. Individuals with less formal education commonly worked as taxi drivers \citep{PakistaniImmigrantsSocial}. Early migration patterns were highly gendered, with typically men coming to the US first to establish themselves financially before sponsoring spouses to reunite with them  or returning to Pakistan to marry and bring a spouse along as dependent (Barratt 2024). These trajectories shaped family formation practices and reinforced traditional gender roles, positioning men as primary economic providers and women as later-arriving caregivers \citep{chaudharyTaleTwoGenerations2024, PakistaniImmigrantsSocial} (Barratt, 2024).

\section{Pakistani Immigrants}

The Pakistani diaspora refers to Pakistani immigrant communities that live transnational lives by maintaining active connections across nations and cultures sustained through languages and diasporic communities’ focus on keeping the intergenerational transmission of culture going \citep{chaudharyTaleTwoGenerations2024, bhatiaTheorizingIdentityTransnational2009}. First generation Pakistani immigrants make a conscious effort to keep their Pakistani cultural identity alive beyond the geographical boarders of their homeland, Pakistan, while also adjusting into the host culture \citep{chaudharyTaleTwoGenerations2024}. This practice challenges children of immigrant parents who only have references, stories and immediate family modeling the traditions and values, which do not exist outside of their homes or diasporic community.

 Existing literature has examined Pakistani immigrant families across a range of contexts, often positioning them within broader ethnic or religious categories, such as South Asian Americans, Asian Americans, or Muslim Americans. Much of this research has approached Pakistani families through lenses of acculturation and assimilation \citep{barkerAcculturationCommunicationFamily2019, bhatiaTheorizingIdentityTransnational2009, kleinRelationshipAcculturationMental2020, myers-wallsLivingLifeTwo2011}. Assimilation refers to the bilateral process through which the characteristics of immigrant groups and host societies start to resemble one another, resulting in a decrease in ethno-racial distinctions \citep{beanAssimilationModelsOld2006, leeAsianAmericanAssimilation2024}. This process, also known as integration or incorporation, encompasses both economic and sociocultural dimensions \citep{leeAsianAmericanAssimilation2024}. While scholars have spent time developing the models, Asian scholars have added a different perspective to the scholarship by critiquing the approach of  ''one size fits all'' adopted by US scholars studying the process of assimilation and acculturation \citet{bhatiaTheorizingIdentityTransnational2009} critique the work on migration being discussed mainly as a series of fixed phases and stages that do not account for the specific culturally distinct and politically entrenched experiences of newer, non-European transnationals \citep{leeAsianAmericanAssimilation2024, bornsteinImmigrationAcculturationParenting2020, bhatiaTheorizingIdentityTransnational2009, beanAssimilationModelsOld2006, ngAreMuslimImmigrants2022, samAcculturationConceptualBackground2006}. \citet{leeAsianAmericanAssimilation2024} point out that the racial experience of Asian Americans challenges traditional theories of racial disadvantage and assimilation. Despite facing xenophobia, racism, and anti-Asian violence, Pakistani immigrants are seen to be relatively successful, leading to an assimilation paradox. Primarily, these studies are sporadic and in the disciplines of psychology, sociology, psychiatry, and educational studies. Prior scholarship has focused on topics, such as parental involvement in children’s academic performance and schooling, as well as broader questions of immigrant identity formation within acculturative contexts. While this body of work offers important insights, it often treats Pakistani immigrants as part of aggregated groups, limiting attention to their distinct family communication practices and culturally specific identity negotiations \citep{chaudharyTaleTwoGenerations2024, leoHighExpectationsCautionary2022, iqbalGenerationGameParenting2018, mohammad-arifMasalaIdentityYoung2000, rasoolBangladeshiIndianPakistani2020, kaliaBufferingParentchildConflict2022, ngAreMuslimImmigrants2022, chanEarlyLookParental2020, attias-donfutIntergenerationalRelationshipsMigrant2017, chungGenderEthnicityAcculturation2001, kleinRelationshipAcculturationMental2020, deepakParentingProcessMigration2005, leeAsianAmericanAssimilation2024, mohamedReligionAsianAmericans2023, masoodGenderFamilyCommunity2009, sirinHyphenatedSelvesMuslim2007, mohammad-arifMasalaIdentityYoung2000, maeshimaMentalHealthStigma2022}.
 
 However, over time, a few studies have specifically focused on the Pakistani American population in context of their ethnic and religious identity, acculturation, assimilation, mental health therapy and counseling, challenges, and gender roles \citep{aghideh2020a, rasoolBangladeshiIndianPakistani2020, bashirUnderstandingContributingFactors2018, chattooYoungPeoplePakistani2004, khalequeDifferentialAcculturationPakistani2015, ghaniPakistaniMigrantsUnited2016, chaudharyTIESTHATBIND2021, thornhillIntersectionalFeministApproach2024}. These approaches are important in addressing the need for research that centers the Pakistani American family as a distinct group and examines how their everyday interpersonal communication practices, storytelling, and parenting values shape identity and meaning making across generations. 
 
Studies have shown children who attend school are often the first in their families to become fluent in the language and culture of the host country \citep{liuIdentityBiculturalMulticultural2017, shinCommunicationTheoryIdentity2017}. Due to this competency, immigrant parents tend to rely excessively on their children for emotional support and instrumental assistance, leading to a role reversal where the child acts as an agent for their parents in formal situations. As a result, these children become adult-like parental figures for their families. During this process of parentification, immigrant children's enacted identity and relational identity are changed based on their roles in the family. Their relationships with parents and siblings define their relational identity within the family unit \citep{attias-donfutIntergenerationalRelationshipsMigrant2017, shinCommunicationTheoryIdentity2017}.

At the same time, parents want to keep the familial hierarchy at home where children remain dependent on them until they are married. Emerging adults may increasingly endorse norms that are more consistent with American culture than with South Asian culture, which their parents expect them to adhere to. Previous research has also shown that gender norms, dating, and marriage are domains of significant conflict in South Asian families \citep{kaliaBufferingParentchildConflict2022, bornsteinImmigrationAcculturationParenting2020, kaliaWhenMisterRight2009}. Existing literature shows familial relationships of South Asian Americans and sociopolitical circumstances impact their worldview and well-being \citep{kaliaBufferingParentchildConflict2022, masoodGenderFamilyCommunity2009}.

Until the 1990s, scholarly attention concentrated on first-generation immigrants (people who are born in different country than the one they immigrated to as adults who intend to settle outside of their homeland or country of origin) . The scholarship was redirected from adult immigrants to the children of immigrants, also called second-generation immigrants (children born to first-generation immigrants in the country of immigration). Soon in the late 20th century in the US, and soon, second-generation immigrants grew into a significant area of intellectual inquiry in migration studies \citep{zhouGrowingAmericanChallenge1997, zhouDivergentDestiniesChildren2019}.In the recent past the focus of immigrants scholarship has started looking at the variants of the second-generation of immigrants counting them as 1.5 generation of immigrants, young children who immigrated with their parents to another country from the country they are born in \citep{parkLanguageIdentitySense2026}  Immigration scholars, sociologists, lawmakers, and policymakers have aimed to understand how immigrants assimilate into the host society and integrate into its political system and institutions. Studies also looked at the influence of Islam, cultural and gender norms of the country of origin, as well as children’s attitudes towards gender roles \cite{bhatiaTheorizingIdentityTransnational2009, zhouDivergentDestiniesChildren2019}.

 Post 9/11 scholarship focused more closely on Muslims and Pakistanis' cultural assimilation and acculturation based on their identity as Muslims, gender roles, and differential acculturation among Pakistani immigrant parents and children \citep{ngAreMuslimImmigrants2022, aghideh2020a}. Results showed discrepancies in parent-child acculturation levels leading to intergenerational acculturation gaps. Sometimes, these gaps led to family conflict, children’s alienation, conduct problems, alcohol and substance abuse, psychological maladjustment, and depression \citep{khalequeDifferentialAcculturationPakistani2015}. Further studies have explored immigrant acculturation experiences in the context of bicultural identity. The findings showed that scholars find it hard to explore the holistic bicultural identity; researchers from different disciplines typically focus on a single aspect, such as language or values \citep{liuIdentityBiculturalMulticultural2017}.

\subsection{Pakistani Family Culture}

Pakistani immigrant families are characterized by their strong commitments towards familial and cultural preservation rooted in collectivist practices while navigating the complexities of the contrasting American individualistic culture \citep{PakistaniImmigrantsSocial, chaudharyTaleTwoGenerations2024}. Family norms such as conformity, respect for parents and elders, and strong emotional bonding with extended family are pivotal, and cohesion is valued. Individuals from collectivistic cultures tend to conform more to family and social norms than personal interests and choices; therein, first-generation Pakistani Americans remain loyal to the standards and customs they learned from their parents. On the other hand, the US culture is individualist, and, thus, prioritizes self-growth, self-reliance, and autonomy, which second-generation Pakistani Americans are exposed to at schools and other social venues leading to an identity gap \citep{saleemFamilyCommunicationProsocial2021}. Identity gaps are ``emotional, behavioral, or cognitive dissonance, which results from experiencing negative communication and relational outcomes'' \citep[p.~225]{hechtCommunicationTheoryIdentity2021}.

The Pakistani family unit subscribes to a rigid and hierarchical structure where every member holds a role with an associated set of responsibilities, which defines their initial personal and ascribed identity. The goal of this structure is to maintain solidarity and harmony both within the family unit and outside, with an emphasis on maintaining the respectable communal identity of the family. They also value the collective needs and desires of the family over the individual. Pakistani parents, thus, tend to exercise higher levels of control and punishments to ensure the assigned roles are followed \citep{anjumRELATIONSHIPPARENTINGSTYLES2019, chanEarlyLookParental2020}.

A core concept in child-rearing in Pakistani Families is \emph{Tarbiyat} an Urdu word that is often translated as ``training'' or ``supervision'' in English but has a much broader meaning in the context of raising children. In practice, \emph{tarbiyat} carries the connotation of guiding your children and training them to be a productive part of society. Parents are also referred to as having the role of a guide or a supervisor in Islamic literature \citep{chanEarlyLookParental2020}. 

Another related concept of family honor, called \emph{izzat}, is central to Pakistani families and holds a significant place within the community. The concept of honor is tightly knitted with conformity, obedience, and academic and financial success. Parenting styles, thus, tend to be more authoritarian and rigorous \citep{kaliaBufferingParentchildConflict2022, raufRelationshipAuthoritarianParenting2017}.

\subsection{Gender Roles}

Conformity is often valued over individual expression and exploration within family life as well as in social and professional contexts. These expectations are deeply gendered, with parents holding different and often more restrictive standards for daughters than for sons. At times, cultural values are maintained through the surveillance of girls' dressing choices, peer relationships, and social behavior. Such monitoring significantly shapes Pakistani immigrant women’s relational roles and influences how they negotiate identity within family and community contexts \citep{chaudharyTaleTwoGenerations2024, felicianoGenderedPathsEducational2005}.

This scrutiny also puts women under a unique set of pressures to fulfill the traditional gender roles, which are ascribed social roles and responsibilities, and characteristics of a gender \citep{PakistaniImmigrantsSocial, chaudharyTaleTwoGenerations2024}. Immigrant parents place higher expectations on daughters, reflecting broader efforts to safeguard family honor and cultural continuity, by making immigrant women bearers of cultural values and communal identity as well as transmitters of culture across generations \citep[p.~3]{chaudharyTaleTwoGenerations2024}.

Pakistani immigrant culture is deeply rooted in family interdependence, positioning everyday family interactions and shared stories as central sites of identity formation for children. In contrast, dominant US cultural norms emphasize individualism, autonomy, and self-expression, placing many immigrant families in a state of ongoing tension as they navigate competing value systems across public and private contexts such as work and school versus home life \citep{chaudharyTaleTwoGenerations2024, saleemFamilyCommunicationProsocial2021, bhatiaTheorizingIdentityTransnational2009}. These cultural contradictions often permeate everyday decision-making and shape how family members negotiate belonging and identity.

Within this framework, gendered expectations play a significant role in structuring family life. Men are commonly positioned as the primary economic providers, while women are expected to assume homemaking and caregiving roles, with participation in paid labor often contingent upon spousal approval. Such expectations can generate familial strain when children and parents encounter conflicting cultural values related to gender roles, independence, and major life decisions, highlighting how cultural transitions are lived, negotiated, and contested within immigrant family communication \citep{kaliaBufferingParentchildConflict2022}.

\subsection{Religion}

Religion is an integral part of Pakistani identity. Being the majority religion, Islam is deeply intertwined with cultural traditions and values. Religion continues to play a vital role in the construction of transnational identity and community for Pakistani immigrants in North America \citep{afzalPakistaniFamilies2021}. The majority of Muslim Pakistani American engages with local mosques and practice religious rituals \citep{ghaniPakistaniMigrantsUnited2016}. Studies show that for many Pakistani Americans religious and cultural heritages are intertwined  \citep{bashirUnderstandingContributingFactors2018}. 

Pakistani Americans are often defined based on dual identity of Muslim Americans, which became a source of great stress and threats post 9/11. Western media projected Muslims Americans as an extremist religious group with links to terrorists who ``hate and kill Americans'' \citep{bilalPortrayalPakistanSilver2021}. Since October 7, 2023, with the escalation in the Gaza conflict, there has been another wave of hatred and xenophobia against Muslim Americans \citep{yancey-braggAntiMuslimDiscriminationReached2025, deroseAntiMuslimBiasReports2024, baileyCAIRReceivedUnprecedented2023}. The western media in the past have shown hatred and violence against Muslim Americans, including women and children \citep{deroseAntiMuslimBiasReports2024, singhAntiMuslimIncidentsJump2024}.

One such study on hyphenated identity with young Muslim Americans in post 9/11 American shows how identities are experienced in nuanced ways \citep{chaudharyTaleTwoGenerations2024, sirinHyphenatedSelvesMuslim2007}. The authors found that boys and girls experienced their Muslim American identities in different ways. Young men often saw their ``Muslim'' and ``American'' identities as separate and sometimes even conflicting. In contrast, young women tended to experience these identities as more connected and overlapping, which created both challenges and opportunities as they moved between family, community, and wider American society \citep{sirinHyphenatedSelvesMuslim2007}. This research shows that identity navigation is shaped by both gender and social context.

\section{Retrospective Storytelling and Communication of Identity}

I adopted the Communication Theory of Identity (CTI) and Communicated Narrative Sense-Making (CNSM) as complementary frameworks because together they allow for a more holistic understanding of identity as both layered and processual.

CTI provides the foundation for this study by conceptualizing identity as interpenetrating and simultaneously active layers. CTI framework is particularly well-suited for examining immigrant families, as prior research has extensively employed CTI to understand how identity is continuously negotiated across cultural, relational, and environmental contexts. The recent inclusion of the fifth frame, the material layer, further strengthens CTI’s applicability, allowing for a more nuanced analysis of how physical, health, territorial and spatial realities shape identity formation. 

By integrating CNSM, this study conceptualizes storytelling not merely as expression, but as a mechanism that activates and connects CTI’s identity layers. Retrospective storytelling becomes the site where identities are interpreted, negotiated, and passed across generations. Together, the framework allowed me to show that retrospective storytelling is the mechanism through which CTI’s identity layers are activated, negotiated, and transmitted across generations.

This study engages with CTI and CNSM simultaneously to demonstrate that storytelling acts as a communication medium and structuring force that actively constructs, transmits, and negotiates identity across personal, enacted, relational, communal, and material layers.

Within households, daily interpersonal conversations take place in the form of stories. Everyday stories function as a central communicative practice through embedded daily interactions. These everyday stories act as a building block, which over time plays a role in shaping individuals’ perception about themselves, cultural norms, religious values, migration experiences, and expectations around gender roles, education, career, marriage, and belonging are shared and interpreted.

Storytelling establishes the foundation of one’s identity in the world; family stories create an individual’s self-perception. Moreover, esteem, familial, relational and social bonds are established. Stories connect parents, children and socialize other members of family and society, help families cope with difficulties, trauma, and stress while communicating family identity and esteem \citep{koenigkellasFamilyStoriesStorytelling2012}. At our core, we are the stories of our family.

Parents play a crucial role in shaping a child's worldview as they act as the socializing agents bridging the micro and macro contexts of their child's life (Kalia et al., 2022). This bridging is built on storytelling as a means of connecting and socializing (Koenig Kellas 2022).  Thus, the way the parents regard the outside world can have a major effect on how the children react to it later.  (Kalia et al., 2022). Parents help their children adopt social norms and expectations from the broader environment. Storytelling operates alongside parenting as a powerful socializing force. Family stories do not merely recount experiences; they actively shape how family members interpret the world, understand themselves, and relate to one another.

Storytelling serves multiple functions in creating identity and socializing, managing stress, and fostering connection through shared values and beliefs, thereby, reinforcing family identity and relational bonds \citep{koenigkellasCommunicatedNarrativeSenseMaking2022, koenigkellasFamilyStoriesStorytelling2012}. Through stories about migration, sacrifice, success, and struggle, immigrant parents communicate expectations, moral lessons, and cultural priorities, while children learn how to locate themselves within family and community narratives.

\subsection{Communicated Narrative Sense Making Theory (CNSM)}

CNSM shifts attention from what family roles are to how the roles are communicated. Family members are talked into being who they as well as how they make sense of their values and beliefs and transfer them to others. The insights of identity development, negotiation and gaps emerge, widen, and are repaired through everyday stories of migration, parenting, sacrifice, and belonging.

Viewed through the lens of CNSM, these family stories have a long-lasting effect on both the storyteller and the listener as families pass down intergenerational values and meaning making. CNSM theory focuses on understanding how the process and content of storytelling serve the ``functions of creating identity, socializing, coping with difficulty, and connecting with other people'' \citep[p.~118]{koenigkellasCommunicatedNarrativeSenseMaking2022}. CNSM theory uses these guiding heuristics: retrospective storytelling, interactional storytelling, and translational storytelling. Retrospective storytelling focuses on the content and the structure of storytelling, exploring its impact on both the storyteller and the listener \citep{koenigkellasCommunicatedNarrativeSenseMaking2022, koenigkellasFamilyStoriesStorytelling2012}. 

The Retrospective Storytelling heuristic in CNSM theory examines how individuals frame the stories they tell and hear \citep[p.~120]{koenigkellasCommunicatedNarrativeSenseMaking2022}. The Interactional storytelling heuristic explicitly focuses on the communicative processes that characterize the collaborative storytelling process. It examines verbal and nonverbal storytelling behaviors of all involved in the interactional process and its link to individual and relational well-being. The translational storytelling heuristic uses narrative theories, methodologies, and empirical findings to develop interventions to enhance individual and collective well-being. This approach emphasizes the synthesis of knowledge. Interventions are crafted by synthesizing the insights derived from research within the retrospective and interactional storytelling heuristics and draws from the extensive interdisciplinary studies in psychology, narrative identity, and storytelling research \citep{koenigkellasCommunicatedNarrativeSenseMaking2022}.

Koenig Kellas recommends the CNSM for the scholars studying ``the story content as the object of their investigation and examining the individual, relational, and/or intergenerational impact of stories and storytelling can and should ground their studies in CNSM theory'' \citep[p.~119]{koenigkellasCommunicatedNarrativeSenseMaking2022}. I focus on the process and the content of the stories from the past and its role in defining the identities across generations. I utilized the retrospective storytelling heuristic of CNSM in this project along with CTI.

\subsection{Retrospective Storytelling}

The retrospective storytelling heuristic of CNSM theory examines the way people frame the stories they tell and stories they hear and the highly significant and meaningful link between this process of storytelling with individual and relational identity, meaning-making, and wellbeing \citep[p.~120]{koenigkellasCommunicatedNarrativeSenseMaking2022}. Hence, the first two propositions of CNSM fall under the retrospective storytelling heuristic explained below:


\textbf{Proposition 1}: The content of storytelling reveals individual, relational, and intergenerational meaning-making, values, and beliefs. 

\textbf{Proposition 2}: Retrospective storytelling content that is framed positively (e.g., redemptively, prosocially, and affectively positive.) This oral sharing must be characterized by high levels of agency and coherency. Then, the narrating will be positively related to individual and relational health and wellbeing \citep[p.~120]{koenigkellasCommunicatedNarrativeSenseMaking2022}.

CNSM theory recognizes and identifies the significant role storytelling and narratives play in defining an individual’s identity, along with the social standing and power dynamics that can be determined by analyzing the verbal and nonverbal content, processes, or structures of stories.

CNSM theory roots itself in Interactional Sense-Making (ISM), meaning a set of behaviors, which conceptualizes meaning as something that is co-constructed through interaction rather than residing solely within individuals. ISM focuses on the process of communication: how meaning emerges through turn-taking, verbal responses, nonverbal cues, pauses, tone, and embodied behaviors as people interact with one another.

Building on ISM, CNSM attends specifically to how families make sense of experiences through storytelling as an interactional practice. This means that researchers do not analyze stories as finished narratives, but as communicative events that unfold in real time and are shaped by who is speaking, who is listening, and how participants respond to one another. The communicated perspective-taking (CPT) refers to how individuals explicitly or implicitly signal understanding, validation, disagreement, or reinterpretation of another person’s perspective during interactions. In family storytelling contexts, CPT reveals how family members align with, resist, or negotiate each other’s interpretations of past events \citep{koenigkellasCommunicatedNarrativeSenseMaking2022}.

CNSM offers a valuable perspective on family stories and storytelling and can provide a lens for understanding the steps toward identity building. The link between storytelling and health is founded on the fundamental roles of narratives and storytelling. People tell stories to establish and shape identity as well as to  educate and socialize others on cultural, societal, and personal values, morals, beliefs, behaviors, and actions, which helps them deal with challenges, complexity,  trauma, and uncertainty \citep{koenigkellasCommunicatedNarrativeSenseMaking2022, koenigkellasCommunicatedNarrativeSenseMaking2018, koenigkellasFamilyStoriesStorytelling2012}.

This inquiry derives from CNSM’s retrospective storytelling heuristics, which also served as the guiding theoretical lens. I propose the first two research questions from the first two propositions of the CNSM. Based on the propositions of retrospective storytelling, here are the first two research questions of my inquiry 

\begin{enumerate}
    \item How do Pakistani American families employ retrospective storytelling and narrative sharing to make sense of their identity within the household and in public?
    \item How are the content and process of retrospective storytelling (individual, relational, and intergenerational meaning-making of values and beliefs) reflected in Pakistani American immigrants’ material frame of identity?
\end{enumerate}

Family stories offer a glimpse into the culture of a family to outsiders \citep{koenigkellasFamilyStoriesStorytelling2012}. Families inherit stories that shape a sense of family identity; family stories also help individuals to understand, express, and construct their identities. These narratives are influenced by families' perceptions of their identities and are one of the primary ways families reflect themselves. The sense of family identity perpetuated through these stories carries family traditions across generations. Retrospective storytelling is a powerful tool that allows families to pass on their values, norms, and cultures from generation to generation. This type of storytelling involves sharing stories and experiences passed down through the family, often across different countries and cultures. These stories are significant as they help shape the next generation's identity. 

Immigrant parents are torn between the challenges of meeting the family expectations and values they were born into and attempting to fit into the culture their children are born in. The transfer of family values of previous generations to their children through storytelling at times clashes with the lived experiences of children, creating confusion and identity gaps \citep{barclayFamilyMemoryIdentity2021, koenigkellasCommunicatedNarrativeSenseMaking2022, koenigkellasFamilyStoriesStorytelling2012}.

Everyday interactions between family members, including conflicts, can become border disputes where cultural differences are given voice and identities are reconstructed. Prior research studying South Asian families found that family conflict increased with the number of years the family spent in the US. Specifically, participants experienced more conflicts related to education and career goals \citep{onayLivingClashSecular2023, kaliaBufferingParentchildConflict2022, giguereLivingCrossroadsCultural2010}.

This dissertation extends existing work examining how Pakistani immigrant families narratively negotiate these cultural expectations through retrospective storytelling, focusing on how parents and children interpret, justify, resist, and rework identity within both private (home) and public (school, work, community) contexts. Drawing on the Communication Theory of Identity (CTI) and Communicated Narrative Sense-Making (CNSM), this study shifts attention from what family roles are to how they are talked into being \citep{duckTalkingRelationshipsBeing1995} and how identity gaps emerge, widen, or are repaired through everyday stories of migration, parenting, sacrifice, and belonging.

By centering lived narratives, this project foregrounds the emotional, embodied, and relational dimensions of identity negotiation that remain underexplored in existing South Asian family communication research. Pakistani Americans remain a marginalized and understudied community that experiences intersectional discrimination based on xenophobia and racism in the US \citep{aghideh2020a}. Very few public images in the US show the cultural values, strengths, and family practices of Pakistanis. Other times, the Pakistani population falls under the shadow of Indian culture and values \citep{bashirUnderstandingContributingFactors2018}.


\section{Identity}

Identity is a complex phenomenon, but it can be seen as a sense of who we are and where we come from. The act of storytelling and the memories shared within families play a central role in establishing, negotiating, understanding, and adjusting personal identity \citep{barclayFamilyMemoryIdentity2021}. Identity negotiation refers to the process of managing and reconciling different aspects of one's self-concept in response to social interactions, life events, and personal experiences \citep{hecht2002ResearchOdyssey1993, jungElaboratingCommunicationTheory2004, kuiperBridgingGapsAdvancing2023, shinCommunicationTheoryIdentity2017}. This process involves navigating the complexities of identity formation, maintenance, and adaptation, both internally and externally. Individuals engage in this ongoing process as they strive to understand and express their identities in a changing world \citep{kuiperCommunicationTheoryIdentity2021, shinCommunicationTheoryIdentity2017}. Identity negotiation is a set of complex communicative processes in which people aim to balance their social interactions and their identity. This task involves an agency to attain social and psychological coherence. However, engaging in identity negotiation can put individuals at risk of alienation, estrangement, and hostility, especially those with stigmatized identities. This complex process requires a lot of communication and effort to achieve the desired outcome \citep{trinhUsingCommunicationTheory2023}.

\subsection{Communication Theory of Identity}

Michael Hecht proposed the communication theory of identity (CTI) in 1993 based upon both post-positivistic and interpretive frameworks \citep{hecht2002ResearchOdyssey1993, shinCommunicationTheoryIdentity2017}. CTI draws from social identity theory and suggests that identity contains multiple frames, each representing a different aspect of one’s identity. All the frames overlap with each other to communicate an individual's complete identity. CTI focuses on the mutual influences on both communication and identity, positing that communication is identity and that identity is a communication field \citep{hechtCommunicationTheoryIdentity2021, jungElaboratingCommunicationTheory2004, shinCommunicationTheoryIdentity2017}. According to \citet{hechtCommunicationTheoryIdentity2021}, CTI draws upon \citet{tajfel1979integrative}'s articulation of the postpositivist inter-group approach. Henri Tajfel and John Turner worked together on the concepts of shared social identity, self-categorization, and categorization of others based on cognitive and social reactions to shared attributes between groups \citep{hechtCommunicationTheoryIdentity2021, hecht2002ResearchOdyssey1993}. 

The CTI raises eight major premises, namely: Identities have individual, social, and communal properties; identities are both enduring and changing; identities are affective, cognitive, behavioral, and spiritual; identities embody both content and relationship levels of interpretation; identities involve both subjective and ascribed meanings; identities are codes expressed in conversations and define membership in communities; identities’ semantic properties manifest in core symbols, meanings, and labels; and identities prescribe modes of appropriate and effective communication \citep[p.~79]{hecht2002ResearchOdyssey1993}.

CTI captures the dynamic and fluid nature of identity through the multi-layered frames and the interpenetration of various layers. Originally, \citet{hecht2002ResearchOdyssey1993} proposed the four layers of identity; personal, relational, enacted, and communal. \citet{kuiperCommunicationTheoryIdentity2021} introduced ``material'' as the fifth frame of identity, which refers to the physical aspects of self, health, and the environment.

Personal identity refers to the concept of self, how an individual defines and views themselves. As social entities, we are given specific meanings of self by others that influence our motivations. These internalized meanings of self are projected in actions and expectations of the fulfillment of those ascribed identities \citep{kuiperCommunicationTheoryIdentity2021, jungElaboratingCommunicationTheory2004}. This layer acknowledges that identity is not a fixed concept but is instead shaped by various factors along with one’s self-perceptions \citep{hechtCommunicationTheoryIdentity2021, shinCommunicationTheoryIdentity2017}.

Enacted identity is defined as identity performed in front of others. The act of communicating with others enacts identity \citep{jungElaboratingCommunicationTheory2004, hecht2002ResearchOdyssey1993}.

Hence, the performance of identity influences communication. The communication of our sense of enacted identity is reinforced by external social interaction \citep{shinCommunicationTheoryIdentity2017}. The concept of the enacted identity frame is based on \citet{goffmanPresentationSelfEveryday2007} idea that the world is like a stage where we all play different roles based on our perception of what others expect from us \citep{jungElaboratingCommunicationTheory2004}. This frame is mainly defined by communication and social behaviors and emphasizes two things about expression. Firstly, identities do not just appear suddenly but develop gradually and keep evolving over time. Secondly, identity is formed and acted out through our behaviors, the roles we choose to play, and the symbols we adopt. This frame also suggests that identity is always present in our communication, and we may not always intentionally communicate our identity messages \citep{kuiperCommunicationTheoryIdentity2021, shinCommunicationTheoryIdentity2017}.

The relational identity frames address who we are in terms of others. In other words, the perception of ``how others view who I am'' creates a sense of ascribed relational identity. Relational identity contains four sublevels. First, individuals’ self-identification is based on internalizing others’ views about themselves – ascribed relational identity. Second, an individual’s identity arises partly in the context of their relationships with others: someone’s spouse, friend, mother, or daughter.  Thirdly, one’s identity can coexist with their other identities (identities in personal, social, and professional lives). Fourth, the approach to the relational layer entails seeing the relationship itself as a unit of identity \citep{jungElaboratingCommunicationTheory2004}. Relational identity explains that identity is not solely self-defined; it is also constructed through interactions with others and our roles in various social contexts \citep{hechtCommunicationTheoryIdentity2021}. 

The communal layer articulates how one’s identity is defined as a part of a larger group or community. This layer defines identities within collectives. The communal layer is a group characteristic rather than individualistic \citep{jungElaboratingCommunicationTheory2004}. Media representations shape communal identities based on religion, ethnicity, education, politics, or other shared characteristics. 

CTI's four layers provide a comprehensive framework that captures the complexity of identity by considering both individual and social aspects. All layers involve communication at the intrapersonal and interpersonal levels as individuals express their multilayered identities. All the layers interpenetrate with each other to some extent and influence the performance of the identity. While CTI conceptualizes identity as operating across personal, enacted, relational, and communal layers, more recent scholarship has expanded this framework to include the material frame, emphasizing the role of physical attributes, health, spatial environments, and territorial location in shaping identity processes. The material frame draws attention to how identities are not only communicated but also embodied, situated, and interpreted through visible and structural conditions \citep{kuiperCommunicationTheoryIdentity2021}.

In 2021, Kim Kuiper proposed material as the fifth frame of CTI. She argues that ``CTI has roots within the intercultural communication emphasis, and this current theoretical expansion allows for continued scholarship such as the physical distinction of passing or having the appearance of a member of a different culture.'' The material frame includes the following dimensions that shape and influence our identities.: 

\begin{itemize}
    \item Our physical attributes,
    \item Our health factors, and 
    \item Our territory and environment
\end{itemize}

Physical attributes refer to the body's physicality, which generates the first visible impression to the outer world. Physical attributes include height, weight, and other traits like agility and posture. Identities emerge from our physical attributes such as appearance, gender, sexual orientation, and health. These attributes can influence the way we see ourselves and how others perceive us \citep{kuiperCommunicationTheoryIdentity2021}.

It is important to acknowledge that the impact of health conditions on our sense of self takes place based on the internalized and social stigma surrounding various health issues, religion, and ethnicity \citep{kuiperCommunicationTheoryIdentity2021, catalanoSupportSocialcognitiveModel2021}. Additionally, the environments we encounter contribute to shaping our identities. Cultural, social, and geographical influences are all factors that play a role in this shaping process \citep{kuiperCommunicationTheoryIdentity2021}. The material frame expands on \citet{jamesConsciousnessSelf1890} notion of the material self, emphasizing the significance of including physicality in identity exploration. This frame recognizes that our bodies, health status, and surroundings shape our identities. By incorporating the material frame, CTI offers a more comprehensive understanding of identity negotiation and how individuals navigate their sense of self within societal contexts \citep{kuiperCommunicationTheoryIdentity2021}. 

This study addresses a gap in prior studies by examining how the material frame functions not as a separate layer, but as a critical site where the alignment and misalignment of other identity layers become visible and experientially salient.

In this sense, the material frame provides a unique analytical lens: through physical appearance, bodily markers (e.g., dress, hijab, skin color), health, and environmental context (e.g., schools, neighborhoods, workplaces), individuals encounter how their identities are read, constrained, or enabled by others. These material conditions shape not only how identity is enacted, but also how it is negotiated across relational and communal expectations.

I argue that the process of retrospective storytelling is projected through the lens of the material frame of the CTI in first-generation immigrants and is used for identity negotiation by their children. CTI explains how an individual develops a sense of identity by articulating various levels of identity formation processes. These stages include personal perception of self, social interactions and relationships, and a collective sense of self in the community \citep{shinCommunicationTheoryIdentity2017}. The originally proposed CTI suggests four layers of identity: personal, enacted, relational, and communal. In recent years, a communication scholar, \citet{kuiperCommunicationTheoryIdentity2021}, proposed a fifth frame, material, that encompasses physical attributes, health factors, and the territories we inhabit. These factors affect the shaping of identities. Health factors, including physical and mental well-being, also impact our identity \citep{kuiperCommunicationTheoryIdentity2021}. Thus, I offer my third research question:

\begin{enumerate}
    \item[3.] How does the material frame of CTI project process identity negotiation between first-generation, one-and-a-half-generation and second-generation immigrants?
\end{enumerate}

\subsection{Identity Gaps}

CTI conceptualizes these layers of identities as interdependent. CTI further adds to the depth of the understanding of identity layers and their interdependence by using the construct of ``interpenetration'' to reflect the process of blending or connecting the four layers of identities. The disconnect between the layers of identities may cause an ``identity gap.'' ``Identity gaps'' happen when individuals experience emotional, behavioral, or cognitive dissonance, which result in them experiencing negative communication and relational outcomes'' \citep[p.~225]{hecht2002ResearchOdyssey1993}. The inconsistencies within the layers of identity lead to conflicted feelings. These inconsistencies cause changes in behavior or relationships as one tries to reduce the identity gap. ``Identity gaps have created perhaps the most prolific line of CTI research'' \citep[p.~226]{hecht2002ResearchOdyssey1993}. 

An identity gap refers to conflict between different layers of one's identity. This tension occurs when an individual sees themselves (their self-perception) differently from how others perceive them. Identity gaps arise from internal conflicts, external pressures, or changes in life circumstances, leading to a sense of fragmentation or dissonance within one's sense of self. An identity gap with respect to the material frame indicates an instance where a discrepancy arises between an individual's physical experiences, such as health conditions or environmental influences, and their self-perception or societal expectations.  

In its application, CTI treats the five layers as interdependent and examines the interpenetration of more than one layer at once to make sense of the juxtaposition of identity layers. The five layers of identity may not always be in sync. Instead, they can be in conflict or mutually exclusive, creating dialectical tension. Despite the contradiction, these frames still coexist and function as integral parts of one’s identity. The disconnect between the layers causes an identity gap. An identity gap can prompt the person to change their behavior or relationships to minimize or resolve discomfort. This response could involve altering the way they communicate, adjusting their expectations, or seeking out new experiences that help them better align with their sense of self. Simply, the inconsistency between the layers of identity is called an identity gap \citep{hechtCommunicationTheoryIdentity2021, jungElaboratingCommunicationTheory2004, kuiperBridgingGapsAdvancing2023, shinCommunicationTheoryIdentity2017}. 

Immigrants to the US make significant adjustments to their identities across various dimensions; transformation is shaped by territory as well as social and physical environments \citep{kuiperCommunicationTheoryIdentity2021}. Moreover, immigrants may find themselves labeled as ethnic minorities within the host culture, a categorization that adds a new dimension to their communal identity, especially if they originated from relatively homogeneous societies \citep{shinCommunicationTheoryIdentity2017}. This communal identity, influenced by societal perceptions and categorizations, further shapes immigrants' personal identity. 

Additionally, it is essential to consider the material frame and the influence of physical, spiritual, and health factors on immigrant identity. These aspects contribute to the multifaceted nature of identity and highlight the need to explore immigrant experiences beyond cultural adaptation alone. The intersection of media representations and immigrant experiences within the material frame of CTI explores how media manipulates the way we see certain groups of people. For example, the media often portrays men wearing \emph{shalwar kameez}, the national dress of Pakistan, as ``bad guys'' or terrorists, which can negatively affect how Pakistani Muslims are perceived. The way a particular cultural identity is projected through a material frame may challenge one's perception of self at communal, enacted, and personal layers. 

In Pakistani culture, norms surrounding the relationship between children and their parents are determined not only by gender and socio-economic background, but also by cultural values that underlie intergenerational obligations, rights, responsibilities, and familial duty. These obligations are transferred from one generation to another through storytelling. This study investigates how parents and their children navigate the complexity of living in two different cultures through narrative sharing. Exploring the impact of the material frame in identity negotiation of Pakistani American families helps us better understand how identity gaps are bridged for marginalized communities. Therein, the Communicated Narrative Sense Making (CNSM) theory illuminates how the sense-making ability of storytelling illuminates these experiences for different generations.  That lead us to the last research question of this study;

\begin{enumerate}
    \item[4. ] Are there gaps in identity for Pakistani immigrants? If so, what are the gaps and how do these gaps influence identity negotiation?
\end{enumerate}

Identity gaps are additionally significant as this subgroup overwhelmingly consists of early-career professionals who are also in the process of starting families. According to statistics from the Organization for Economic Cooperation and Development on foreign medical doctors in OECD member nations, Pakistan moved up to the second spot in 2023, making it the second biggest source of young foreign doctors in the US \citep{mullanMetricsPhysicianBrain2005, haqPakistanSecondBiggest2023}.

\section{Addressing the Gaps and the Current Study}

This study integrates CTI and CNSM to demonstrate that storytelling is not merely a reflection of identity but a structuring force that actively constructs, transmits, and negotiates identity across personal, relational, enacted, communal, and material layers. By marrying CTI and CNSM theories together, this study shows that identity is shaped and governed by the narratives where storytelling structures are how identities are formed, felt, and carried across the communities and generations. 

Recent studies boost the critical angles within interpersonal and family communication studies (IFC) and the advances of diversity, equity, inclusion, and access in the larger field of communication studies. However, in CIFC, there are several areas that still require further exploration. For instance, more research is needed on the experiences of marginalized communities, such as immigrant families and families with disabilities \citep{faulknerSettingAgendaPoetic2023, manningHammerCriticalTurn2021, mooreWhatCountsCritical2019}. Further, scholars need to study the impact of cultural differences on family communication and relationships, as well as the ways in which families navigate cultural barriers and their differences. 

Critical interpersonal and family communication (CIFC) scholars contend that we need research on the experiences of marginalized communities, such as LGBTQ+ families, immigrant families, and families with disabilities \citep{faulknerSettingAgendaPoetic2023, manningHammerCriticalTurn2021, mooreWhatCountsCritical2019}. These scholars further call for study on the impact of cultural differences on family communication and relationships as well as how families navigate cultural barriers and negotiate their differences \citep{manningHammerCriticalTurn2021}. They suggest that interpersonal and family communication (IFC) would benefit from exploring how theories manifest in different cultural contexts. Critical theories, perspectives, approaches, and traditions could be useful to decolonize the IFC canon, which historically has been shaped by Western-centric and often static views of culture. Thus, this study centers the lived experiences and gendered voices of diasporic, non-western, people of color. It provides an insider's perspective to demonstrate how Pakistani American families’ interpersonal and family communication is present in their storytelling. I provide a holistic view through my investigation into the process of identity development and negotiation from two distinct perspectives: First-generation Pakistani immigrants (in the role of parents) and 1.5 generation and second-generation Pakistani immigrants (as children who either accompanied their immigrant parents as their dependents or were born in the US). Specifically, this study examines the role of retrospective storytelling in shaping the worldview and identities of second-generation immigrants given the clash of cultures and values they face \citep{koenigkellasCommunicatedNarrativeSenseMaking2022, koenigkellasFamilyStoriesStorytelling2012}.

The study explores the communication patterns within the Pakistani immigrant households and looks at the process of identity development and negotiation of first-generation immigrants as parents and one-and-a-half and second-generation immigrants (the children of immigrant parents who either accompany their parents at a very young age or are born in the US). I focus on how the storytelling process in these homes empowers parents to share their values and things important to the whole family. The communication process between parents and child works as a building block for establishing the foundation of any child’s identity development; the important factors to study here are the impact of parents' immigration to a completely new culture and the task of identity development and negotiation for themselves along with their children in the alien culture of the new country (i.e.  the US). My overarching question is: How does the process of identity negotiation influence the relationship between parents and children? 

Based on the culture, values, and lived experiences in Pakistan, the worldview of first-generation immigrant parents differs from that of their children who grow up in a different culture and worldview \citep{anjumRELATIONSHIPPARENTINGSTYLES2019, chanEarlyLookParental2020, khalequeDifferentialAcculturationPakistani2015}. The religious and cultural heritages are inseparable for first-generation Pakistani Americans \citep{bashirUnderstandingContributingFactors2018}. Previous studies indicate that it is difficult to differentiate whether Pakistani cultural values or religious beliefs have influenced their behaviors. Second-generation Pakistani Americans identify more with American mainstream culture. However, they also maintain their Pakistani heritage by participating in traditional events and upholding the traditional values expected by their parents \citep{bashirUnderstandingContributingFactors2018}. The relationship between children and their parents is influenced by gender, socio-economic background, and cultural values that shape intergenerational obligations, rights, responsibilities, and filial duty. Research suggests that young people born to Pakistani immigrants living in Western countries tend to experience overprotection from their families and are perceived to receive the least support for developing the skills they need to face the world independently \citep{chattooYoungPeoplePakistani2004, rasoolBangladeshiIndianPakistani2020}. This study shows techniques and tools for individuals seeking to improve their family communication and identity negotiation, promoting healthier family dynamics.

Historically, the study of family communication has focused on the family as a private entity, separate from external cultural influences. However, the CIFC perspective suggests that the public and private are interconnected and that identity negotiation within a family cannot be understood as an isolated entity, but rather as embedded in sociocultural contexts and subject to political and cultural influences \citep{suterIntroductionCriticalApproaches2016}. In the case of first-generation Pakistani Americans, the influence comes from two distinct and contrasting socio-cultural environments.

In this study, I present the challenges faced by first, 1.5, and second-generation Pakistani American immigrant families by investigating how they negotiated identities both inside and outside of their homes and dealt with difficult conversations with their children at home using the communicated narrative sense-making theory and the communicated theory of identity.

Thus, the purpose of this study was to learn from the lived experiences of first-generation Pakistani immigrants (in the role of parents), and 1.5 and second-generation immigrants (children of first-generation Pakistani immigrants) through their stories. Migration in the later stages of life as parents brings extra layers of challenges and a disconnect between parents and their children when adapting to different cultural and social environments. Many parents struggle with limited English proficiency as well as accent and feel overwhelmed by their children's experiences in American schools and lifestyles. This divide sometimes leads to feelings of alienation and distrust between parents and children. As they attempt to stay connected to their roots, many immigrant parents hold onto old customs from their home country, which can create a disconnect with their children's experience \citep{aghideh2020a, khalequeDifferentialAcculturationPakistani2015}.

To reiterate the research questions (RQ) of this study are as follows:

\section{Research Questions}

\begin{enumerate}
    \item How do Pakistani American families employ retrospective storytelling and narrative sharing to make sense of their identity within the household and in public?
    \item How does the content and the process of storytelling among first-generation immigrant peers impact individual, relational, and intergenerational meaning-making of values and beliefs? 
    \item How is the material frame of identity reflected in the stories of 1.5- and second- generation of Pakistani immigrants? 
    \item Are there gaps in identity for Pakistani immigrants? If so, what are the gaps and how do these gaps influence identity negotiation?
\end{enumerate}

To answer these research questions, I used Arts-Based Research (ABR), which is a qualitative methodology, and specifically, poetic inquiry. Poetic inquiry allowed me to delve into the lived experiences, viewpoints, and identity negotiations of Pakistani American families, including parents (first-generation immigrants) and Pakistani American children born to first-generation Pakistani Americans of age 18 years and above (one-and-a-half and second-generation immigrants).

In summary, the chapter introduces the key scholarship looking at Pakistani American immigrant families, along with the overview of US immigration policies, Pakistani culture, and family communication.  It also introduces the theoretical frameworks guiding this study, the Communication Theory of Identity (CTI) and Communicated Narrative Sense-Making (CNSM) as well as explaining their relevance for examining identity as a relational, retrospective narrative. 

Finally, the chapter identifies critical gaps in existing scholarship, particularly the limited attention to intergenerational, narrative, and materially situated identity negotiation within Pakistani American families, thereby establishing the foundation for the present research.
